\documentclass[11pt,a4paper,landscape]{article}
\usepackage[T1]{fontenc}
\usepackage[utf8]{inputenc}
\usepackage[provide=*,lithuanian]{babel}
\usepackage{lmodern}
\usepackage{amsmath,amssymb}
\usepackage[margin=0pt]{geometry}
\usepackage{tikz}
\pagestyle{empty}
\setlength{\parindent}{0pt}
\setlength{\parskip}{0pt}

% Spausdinti horizontaliame A4 lape, mastelis 100 %.
% Perkirpus per vidurį gaunami du statūs A5 dydžio lapeliai.
\newcommand{\remelis}[2]{%
  \tikz[baseline={([yshift=-0.6ex]current bounding box.center)}]\draw[line width=0.5pt]
    (0,0) rectangle (#1,#2);}

\newcommand{\uzduotis}[2]{%
  \begin{minipage}[t][25mm][t]{\linewidth}
    \textbf{#1.}\enspace #2\par
    \vspace{3mm}
    \remelis{50.4mm}{12mm}\par
  \end{minipage}\par
}

\newcommand{\lapelis}[2]{%
  \begin{minipage}[t][190mm][t]{128.5mm}
    \small
    \textbf{Mokinio kodas}\enspace\remelis{32mm}{8mm}%
    \hfill\textbf{#1}\par
    \vspace{6mm}
    {\centering\normalsize\textbf{Bazinių įgūdžių testas 1}\par}
    \vspace{2mm}
    {\centering\small Į rėmelius įrašykite tik atsakymus.\par}
    \vspace{7mm}
    \normalsize
    #2
  \end{minipage}%
}

\newcommand{\stulpeliai}[2]{%
  \begin{minipage}[t]{60mm}#1\end{minipage}\hfill
  \begin{minipage}[t]{60mm}#2\end{minipage}%
}

\begin{document}
\null
\begin{tikzpicture}[remember picture,overlay]
  \draw[dash pattern=on 2pt off 3pt,line width=0.4pt]
    (current page.north) -- (current page.south);
  \node[anchor=north west,inner sep=0pt]
    at ([xshift=10mm,yshift=-10mm]current page.north west) {%
      \lapelis{Gerumo variantas}{%
        \stulpeliai{%
          \uzduotis{1}{Išskirkite pilnąjį kvadratą:\\\(2x^2+8x-3\)}
          \uzduotis{2}{Išskirkite pilnąjį kvadratą:\\\(-3x^2+6x+1\)}
          \uzduotis{3}{Atskliauskite:\\\((x+1)^2\)}
          \uzduotis{4}{Atskliauskite:\\\((x-2)^2\)}
          \uzduotis{5}{Atskliauskite:\\\((x+1)^3\)}
          \uzduotis{6}{Atskliauskite:\\\((2x-1)^3\)}
        }{%
          \uzduotis{7}{Atskliauskite:\\\((x-1)(x+1)\)}
          \uzduotis{8}{Atskliauskite:\\\((x-2)(x^2+2x+4)\)}
          \uzduotis{9}{Atskliauskite:\\\((x+1)(x^2-x+1)\)}
          \uzduotis{10}{Atskliauskite:\\\((a+ab)^2\)}
          \uzduotis{11}{Atskliauskite:\\\((3x-1)^2\)}
          \uzduotis{12}{Išskaidykite dauginamaisiais:\\\(8x^3-1000\)}
        }%
      }%
    };
  \node[anchor=north west,inner sep=0pt]
    at ([xshift=10mm,yshift=-10mm]current page.north) {%
      \lapelis{Pagarbos variantas}{%
        \stulpeliai{%
          \uzduotis{1}{Išskirkite pilnąjį kvadratą:\\\(2x^2+8x-3\)}
          \uzduotis{2}{Išskirkite pilnąjį kvadratą:\\\(-3x^2+6x+1\)}
          \uzduotis{3}{Atskliauskite:\\\((x+2)^2\)}
          \uzduotis{4}{Atskliauskite:\\\((x-1)^2\)}
          \uzduotis{5}{Atskliauskite:\\\((x+2)^3\)}
          \uzduotis{6}{Atskliauskite:\\\((2x-1)^3\)}
        }{%
          \uzduotis{7}{Atskliauskite:\\\((x-2)(x+2)\)}
          \uzduotis{8}{Atskliauskite:\\\((x-1)(x^2+x+1)\)}
          \uzduotis{9}{Atskliauskite:\\\((x+2)(x^2-2x+4)\)}
          \uzduotis{10}{Atskliauskite:\\\((a+ab)^2\)}
          \uzduotis{11}{Atskliauskite:\\\((2x-1)^2\)}
          \uzduotis{12}{Išskaidykite dauginamaisiais:\\\(8x^3-1000\)}
        }%
      }%
    };
\end{tikzpicture}
\end{document}
